\documentclass[11pt]{article}
\usepackage[T1]{fontenc}
\usepackage{lmodern,amsmath,amssymb,amsthm,booktabs,array,longtable}
\usepackage[margin=1in]{geometry}
\usepackage[colorlinks=true,urlcolor=blue,citecolor=blue,linkcolor=blue]{hyperref}
\newtheorem{theorem}{Theorem}
\newtheorem{corollary}[theorem]{Corollary}
\title{A short proof of Bala's all-column conjecture using Wildon's binomial identity}
\author{Haobo Ma\\\small ChronoAI Pte Ltd, Singapore
\and Wenlin Zhang\\\small National University of Singapore\\\small ChronoAI Pte Ltd, Singapore}
\date{November 16, 2026}
\begin{document}
\maketitle
\begin{abstract}
We give a proof of Peter Bala's all-column conjecture in the final
FORMULA of OEIS A337389, revision 40 (May 5, 2022), by specializing
the binomial identity proved by Mark Wildon on November 27, 2013,
and identified there by Ira Gessel as terminating Kummer summation.
The identity holds with an indeterminate column parameter and includes
the individual column identities previously proved by Bala and
Blanco G\'omez. We also prove a squared-coefficient identity and Gauss
congruences for every integer parameter. We relate the stronger recorded congruences to
their prior proofs, and derive the companion recurrence and asymptotic
formulas, including growth along the main diagonal.
\end{abstract}

\section{The coefficient identity}
Seiichi Manyama's array in OEIS A337389 is specified by a nested radical
and a binomial sum. Peter Bala conjectured a diagonal-coefficient
description on May 3, 2022; the same line records the cases $k=0,4$
as known \cite{array}. Its full coefficient identity is a specialization
of a classical binomial identity: Wildon gave an elementary interpolation
proof in 2013, and Gessel identified it as terminating Kummer summation
in an answer on the same day \cite{wildon,gessel2013}.
We give that short argument and the exact specialization, so the
connection to the later conjecture is explicit.
The column $k=1$ was already settled by Bala in the August 29, 2025
addition to A082758 \cite{columnone}. Blanco G\'omez proved the
columns $k=2,3$ in Section 3.3 of his October 17, 2026 preprint
\cite{blanco}. Theorem~\ref{main} places those identities within one
polynomial identity for all columns; it makes no claim of a new
binomial evaluation.

For a formal series $F$, write $[x^n]F$ for its coefficient of $x^n$.
All square roots below have constant term $1$. Initially work over
$\mathbb Q[k]$, with $k$ an indeterminate, and define
\begin{align}
 \Delta_k(t)&=1-2(k+4)t+(k-4)^2t^2,\nonumber\\
 d_k(t)&=\sqrt{\Delta_k(t)},\nonumber\\
 A_k(t)&=\sqrt{\frac{1+(k-4)t+d_k(t)}{2\Delta_k(t)}},
 &T(n,k)&=[t^n]A_k(t),\label{defA}\\
 \Phi_k(u)&=\frac{(1+(k-2)u+u^2)(1+u)^2}{(1-u)^2}.
 \label{defPhi}
\end{align}
The radicands have constant term $1$ and the denominators are units.
For a series $R$ with $R(0)=1$, the equations $S(0)=1$ and $S^2=R$
determine $[t^n]S$ successively: twice that coefficient equals
$[t^n]R-\sum_{i=1}^{n-1}[t^i]S[t^{n-i}]S$.
Thus these definitions give unique elements of $\mathbb Q[k][[t]]$.

\begin{theorem}\label{main}
For every integer $n\geq0$, the following is an identity in $\mathbb Q[k]$:
\begin{equation}\label{identity}
 T(n,k)=[u^n]\Phi_k(u)^n
 =\sum_{j=0}^n\binom{2n}{2j}\binom{2j}{j}k^{n-j}.
\end{equation}
In particular it holds for every element $k$ of a characteristic-zero
field, including every nonnegative integer column of the source array.
\end{theorem}

\begin{proof}
We first recall Wildon's identity \cite{wildon}:
\begin{equation}\label{kummer}
 [u^h](1-u)^x(1+u)^{2h-x}
   =(-1)^h4^h\binom{(x-1)/2}{h}\qquad(h\geq0).
\end{equation}
Here $\binom{x}{h}=x(x-1)\cdots(x-h+1)/h!$, including value $1$
at $h=0$. Both sides are polynomials of degree at most $h$ in $x$.
For $x=1,3,\ldots,2h-1$, the degree-$2h$ polynomial
$f(u)=(1-u)^x(1+u)^{2h-x}$ satisfies
$u^{2h}f(1/u)=-f(u)$, so its middle coefficient vanishes.
At $x=2h+1$, the left side is
$(-1)^h\sum_{i=0}^h\binom{2h+1}{i}=(-1)^h4^h$.
Thus the two degree-at-most-$h$ polynomials agree at $h+1$ points,
proving \eqref{kummer}; for $h=0$ both sides are $1$.

Now $\Phi_k(u)=(1+u)^2(1+ku/(1-u)^2)$, whence
\begin{equation}\label{transformed}
 [u^n]\Phi_k(u)^n=
 \sum_{m=0}^n\binom nm k^m
 [u^{n-m}]\frac{(1+u)^{2n}}{(1-u)^{2m}}.
\end{equation}
Set $h=n-m$ and $x=-2m$ in \eqref{kummer}.
With $(a)_r=a(a+1)\cdots(a+r-1)$ and $(a)_0=1$, this gives
\[
 \binom nm[u^{n-m}]\frac{(1+u)^{2n}}{(1-u)^{2m}}
 =\binom nm\frac{4^{n-m}(m+1/2)_{n-m}}{(n-m)!}
 =\frac{(2n)!}{(2m)!(n-m)!^2}.
\]
For the final equality, use
$(m+1/2)_{n-m}=(2n)!\,m!/(4^{n-m}n!(2m)!)$.
Replacing $m$ by $n-j$ proves that the diagonal coefficient is the
sum in \eqref{identity}.

It remains to identify its generating function with \eqref{defA}.
Introduce an indeterminate $b$ and set $k=b^2$ and $t=s^2$.
The homomorphism $\mathbb Q[k]\longrightarrow\mathbb Q[b]$ sending
$k$ to $b^2$ is injective, so an identity after this substitution suffices.
Let
\[
 H_b(s)=\big((1-bs)^2-4s^2\big)^{-1/2}
       =\sum_{\ell\geq0}p_\ell(b)s^\ell.
\]
Expanding first in $s^2/(1-bs)^2$ gives
\[
 p_\ell(b)=\sum_{j=0}^{\lfloor\ell/2\rfloor}
        \binom{\ell}{2j}\binom{2j}{j}b^{\ell-2j}.
\]
Consequently the generating series of the sum in \eqref{identity} is
$E(t)=(H_b(s)+H_b(-s))/2$. If $c=k-4$, then
\[
 H_b(s)^2=\frac1{1-2bs+cs^2},\qquad
 H_b(s)H_b(-s)=\frac1{d_k(t)}.
\]
The latter equality follows by squaring and comparing constant terms.
Taking the square of the even part now gives
\[
 E(t)^2=\frac{1+ct}{2\Delta_k(t)}+\frac1{2d_k(t)}
       =\frac{1+ct+d_k(t)}{2\Delta_k(t)}.
\]
Since $E(0)=1$, uniqueness of the normalized square root yields
$E=A_k$, proving \eqref{identity}, also for $n=0$.
\end{proof}

\section{Consequences and the companion array}
\begin{corollary}\label{consequences}
For $n\geq0$,
\[
 T(n,0)=\binom{2n}n,\qquad T(n,4)=\binom{4n}{2n}.
\]
For nonnegative integer $k$, $T(n,k)$ is a nonnegative integer.
The array in OEIS A337419 equals $T(n,-k)$; hence its column $k$
also has the diagonal description $[u^n]\Phi_{-k}(u)^n$.
The main diagonal in A337388 is $T(n,n)=[u^n]\Phi_n(u)^n$.
\end{corollary}
\begin{proof}
At $k=0$, the binomial sum has just its $j=n$ term, with the
polynomial convention $0^0=1$. Nonnegativity and integrality follow
from the same sum. At $k=4$, use $b=2$ above:
$H_2(s)=(1-4s)^{-1/2}$, whose coefficient of $s^{2n}$ is
$\binom{4n}{2n}$. Substituting $-k$ into \eqref{defA} and
\eqref{identity} gives respectively the NAME and first FORMULA of
A337419. Substituting $k=n$ gives the NAME of A337388
\cite{signed,diagonal}. These are consequences of a polynomial identity,
not limits involving division by $k$ or $k-4$.
\end{proof}

The odd part of the same series supplies a nearby array without another
inversion. Define
\begin{equation}\label{odd}
 U(n,k)=\sum_{j=0}^n\binom{2n+1}{2j}\binom{2j}{j}k^{n-j}.
\end{equation}
\begin{corollary}\label{oddcor}
The array in OEIS A337369 satisfies
\begin{equation}\label{oddgf}
 \sum_{n\geq0}U(n,k)t^n
 =\sqrt{\frac{2}{\Delta_k(t)(1+(k-4)t+d_k(t))}}.
\end{equation}
It also has the polynomial integral description
$U(n,k)=(2n+1)\int_0^1T(n,kr^2)\,dr$,
where the integral acts term by term on a polynomial in $r$.
\end{corollary}
\begin{proof}
The series $O(t)=(H_b(s)-H_b(-s))/(2bs)$ has coefficients
$p_{2n+1}(b)/b=U(n,b^2)$; the numerator is divisible by $bs$.
Over $\mathbb Q(k)$ its square is
\[
 O(t)^2=\frac{1+(k-4)t-d_k(t)}{2kt\Delta_k(t)}
       =\frac{2}{\Delta_k(t)(1+(k-4)t+d_k(t))},
\]
because $(1+(k-4)t)^2-\Delta_k(t)=4kt$.
The final expression is a series over $\mathbb Q[k]$, and
the injection into $\mathbb Q(k)$ extends the identity to $k=0$
as well. Its constant term and $O(0)$ are $1$, proving
\eqref{oddgf}, the source's NAME \cite{oddarray}.
For the integral identity compare coefficients of $k^m$ in
\eqref{identity} and \eqref{odd}: their ratio is
$(2n+1)/(2m+1)$.
\end{proof}

\section{Column conjectures and congruences}
The substitutions $k=1,2,3$ in Theorem~\ref{main} give respectively
the diagonal identities in Bala's April 3, May 2 and May 3, 2022
FORMULA blocks of A082758, A337390 and A245926
\cite{columnone,columntwo,columnthree}.
For clarity their rational functions are
\[
 \Phi_1(u)=\frac{(1+u)(1+u^3)}{(1-u)^2},\quad
 \Phi_2(u)=\frac{(1+u^2)(1+u)^2}{(1-u)^2},\quad
 \Phi_3(u)=\frac{(1+u+u^2)(1+u)^2}{(1-u)^2}.
\]
Expanding the numerators and $(1-u)^{-2n}$ also proves the double
and triple sums displayed as equivalent formulations in those entries;
when $n=0$ their value is $1$. The earlier proofs for these columns
are credited above.

For a prime $p$, let $v_p(a)$ be the exponent of $p$ in a nonzero
integer $a$, with $v_p(0)=+\infty$.
\begin{corollary}\label{gauss}
For every integer $k$, prime $p$, and integers $m,r\geq1$,
\[
 T(mp^r,k)\equiv T(mp^{r-1},k)\pmod {p^r}.
\]
\end{corollary}
\begin{proof}
First suppose $k\geq0$. The Laurent series
\[
 L_k(u)=\frac{\Phi_k(u)}u
 =\frac{(1+u)^2}u+k\frac{(1+u)^2}{(1-u)^2}
\]
has nonnegative integer coefficients and smallest exponent $-1$.
Regard its coefficient of $u^j$ as the number of colors for a letter
of weight $j$. The constant term of $L_k(u)^N$ counts length-$N$
colored words of total weight zero. This set is finite, since each
letter has weight at most $N-1$ in such a word.
Let $b_d$ count rotation orbits of primitive weight-zero words of
length $d$. Every length-$N$ word is a repetition of a unique primitive
word whose length divides $N$; its primitive word also has weight zero.
Consequently
\[
 T(N,k)=\sum_{d\mid N}d\,b_d.
\]
For $N=mp^r$, every divisor of $N$ not dividing $N/p$ has
$p$-adic valuation $v_p(N)\geq r$. Subtract the two divisor sums.
For negative $k$, choose $k'\geq0$ congruent to $k$ modulo $p^r$
and use $T(N,k)\in\mathbb Z[k]$ from Theorem~\ref{main}.
\end{proof}
This proves the conditional Gauss-congruence conclusions in all three
column entries, without restricting the prime.

The stronger statements modulo $p^{2r}$ in the same source blocks
already have proofs. Pan and Sun \cite[Theorem 1.1]{pansun} proved
in 2020 that, for $p>3$ and the central trinomial coefficients
$q_j=[u^j](1+u+u^2)^j$,
\[
 v_p(q_{p\ell}-q_\ell)\geq 2+2v_p(\ell).
\]
Since $T(n,1)=q_{2n}$, putting
$\ell=2mp^{r-1}$ proves Bala's $k=1$ supercongruence for $p\geq5$.
This application was also recorded by Epstein in A082758 on
July 18, 2026 \cite{columnone}.

\begin{corollary}[Consequences of Blanco G\'omez's transfer theorem]\label{super}
For $m,r\geq1$,
\[
 T(mp^r,k)\equiv T(mp^{r-1},k)\pmod {p^{2r}}
\]
holds for $k=2$ and every odd prime $p$, and for $k=1,3$ and
every prime $p\geq5$.
\end{corollary}
\begin{proof}
Write $G(u)=\Phi_k(u)=\exp(\sum_{j\geq1}c(j)u^j/j)$.
The three factorizations are
\[
 \Phi_1=\frac{(1-u^2)(1-u^6)}{(1-u)^3(1-u^3)},\qquad
 \Phi_2=\frac{(1-u^2)(1-u^4)}{(1-u)^4},\qquad
 \Phi_3=\frac{(1-u^2)^2(1-u^3)}{(1-u)^5}.
\]
For $G=\prod_{d\mid P}(1-u^d)^{-a_d}$ with integer $a_d$,
$c(j)=\sum_{d\mid\gcd(j,P)}d a_d$, so $c(pj)=c(j)$ if
$p\nmid P$. In particular $c$ satisfies hypothesis (H2) of
Blanco G\'omez's Theorem A, which transfers the congruences to
$[u^N]G(u)^N$. Apply it with $P=6,4,6$, respectively.
These are the already proved rows of his Corollary 3.3
\cite{blanco}, not new supercongruence claims.
\end{proof}

\begin{corollary}\label{squares}
For every $n\geq0$, as an identity in $\mathbb Z[k]$,
\begin{equation}\label{squareid}
 T(n,k)^2=\sum_{j=0}^{2n}(k-4)^{2n-j}
             \binom{2j}{j}^{\!2}\binom{2n+j}{2j}.
\end{equation}
In particular this proves the squared-coefficient conjecture in
Bala's May 2, 2022 block of A337390. It also yields the recorded
square identities of A245926 and A243946 at $k=3,5$
\cite{columnthree,columnfive}.
\end{corollary}
\begin{proof}
Define the Legendre polynomials by
$(1-2xs+s^2)^{-1/2}=\sum_{h\geq0}P_h(x)s^h$.
The series $H_b$ in the proof of Theorem~\ref{main} gives
$T(n,k)=(k-4)^nP_{2n}(\sqrt{k/(k-4)})$ when $k\ne4$.
Write ${}_aF_b$ for the hypergeometric series whose coefficient at
degree $j$ is $\prod_i(\alpha_i)_j/(j!\prod_h(\beta_h)_j)$ for
numerator parameters $\alpha_i$ and denominator parameters $\beta_h$.
The polynomial $P_{2n}(x)$ is
${}_2F_1(-n,n+1/2;1;1-x^2)$: both are polynomial solutions of
$(1-x^2)y''-2xy'+2n(2n+1)y=0$, normalized by $y(1)=1$.
Clausen's identity \cite[16.12.2]{clausen}, with parameters
$-n,n+1/2$, gives
\[
 P_{2n}(x)^2={}_3F_2(-2n,2n+1,1/2;1,1;1-x^2)
 =\sum_{j=0}^{2n}\binom{2j}{j}^{\!2}\binom{2n+j}{2j}
                  \left(\frac{x^2-1}{4}\right)^j.
\]
Substitute $x^2=k/(k-4)$ and multiply by $(k-4)^{2n}$.
Both resulting expressions are polynomials in $k$, which includes $k=4$.
\end{proof}

\section{Recurrence and growth of the original array}
The proof also verifies two formulas recorded before Bala's conjecture.
These formulas are credited in the source to Manyama and
V\'aclav Kot\v e\v sovec, respectively \cite{array}.
Writing $c=k-4$, differentiation of $H_b(s)$ gives
\[
 (1-2bs+cs^2)H_b'(s)=(b-cs)H_b(s).
\]
It follows that $p_0=1$, $p_1=b$, and, for $\ell\geq2$,
\begin{equation}\label{prec}
 \ell p_\ell=(2\ell-1)b p_{\ell-1}-(\ell-1)c p_{\ell-2}.
\end{equation}
To eliminate the odd coefficients for $n\geq2$, apply \eqref{prec}
at $\ell=2n,2n-1,2n-2$. Multiply the first equation by
$(2n-1)(4n-5)$, substitute the second for
$(2n-1)p_{2n-1}$, and the third for
$(4n-5)b p_{2n-3}$. The coefficient of $p_{2n-2}$ is
\[
 (4n-1)(4n-3)(4n-5)k
 -(2n-1)^2(4n-5)c-(4n-1)(2n-2)^2c.
\]
It equals $2(4n-3)(4(k+4)n^2-6(k+4)n+k+6)$.
Thus $T(0,k)=1$, $T(1,k)=k+2$, and
\begin{align}
 n(2n-1)(4n-5)T(n,k)
  ={}&(4n-3)\big(4(k+4)n^2-6(k+4)n+k+6\big)T(n-1,k)
 \nonumber\\[-2pt]
 &{}-(k-4)^2(n-1)(2n-3)(4n-1)T(n-2,k).
 \label{rec}
\end{align}
This is precisely Manyama's August 28, 2020 recurrence in A337389.
The negative-parameter recurrence of A337419 follows by $k\mapsto-k$.

\begin{corollary}\label{oddrec}
For $n\geq2$, the odd companion has $U(0,k)=1$, $U(1,k)=k+6$, and
\begin{align*}
 n(2n+1)(4n-3)U(n,k)
 ={}&(4n-1)\bigl(4(k+4)n^2-2(k+4)n-k-2\bigr)U(n-1,k)\\
 &{}-(k-4)^2(n-1)(2n-1)(4n+1)U(n-2,k).
\end{align*}
This is Manyama's August 29, 2020 recurrence in A337369.
\end{corollary}
\begin{proof}
For any $h\geq4$, elimination of $p_{h-1},p_{h-3}$ from
\eqref{prec} gives
\begin{align*}
 h(h-1)(2h-5)p_h
 ={}&\bigl((2h-1)(2h-3)(2h-5)k\\
 &\quad -(h-1)^2(2h-5)(k-4)
              -(2h-1)(h-2)^2(k-4)\bigr)p_{h-2}\\
 &{}-(2h-1)(h-2)(h-3)(k-4)^2p_{h-4}.
\end{align*}
Set $h=2n+1$, use $p_{2n+1}=bU(n,b^2)$, and cancel $2b$
as a polynomial identity. The middle factor becomes
$(4n-1)(4(k+4)n^2-2(k+4)n-k-2)$ after division by $2$.
\end{proof}

For fixed real $k>0$, put $b=\sqrt k$ and $M=b+2$.
Taking the constant term of $(b+z+z^{-1})^{2n}$ yields the sum
in \eqref{identity}, so
\[
 T(n,k)=\frac1{2\pi}\int_{-\pi}^{\pi}(b+2\cos\theta)^{2n}\,d\theta.
\]
The absolute value of the base has its unique maximum $M$ at $0$.
Near $0$, $(b+2\cos\theta)/M=1-\theta^2/M+O(\theta^4)$.
Outside any fixed neighborhood of $0$ its absolute value is at most
$M(1-\varepsilon)$ for some $\varepsilon>0$.
Inside a sufficiently small neighborhood it is positive and its
normalized $2n$-th power is bounded by $e^{-a n\theta^2}$ for a
constant $a>0$. Set $y=\sqrt{2n/M}\,\theta$ there and use dominated
convergence. The limiting integral is
$\int_{\mathbb R}e^{-y^2}\,dy=\sqrt\pi$, giving
\begin{equation}\label{asympt}
 T(n,k)\sim\frac{(2+\sqrt k)^{2n+1/2}}{\sqrt{8\pi n}}
 \qquad(n\longrightarrow\infty,\ k>0\text{ fixed}).
\end{equation}
This proves Kot\v e\v sovec's August 31, 2020 formula.

\begin{corollary}\label{growth}
For each fixed real $k>0$,
\[
 U(n,k)\sim\frac{(2+\sqrt k)^{2n+3/2}}{\sqrt{8k\pi n}}.
\]
For any real sequence $b_n\geq2$ with $b_n=o(n)$,
\[
 T(n,b_n^2)\sim\frac{(b_n+2)^{2n+1/2}}{\sqrt{8\pi n}}.
\]
In particular the two displayed equivalents for the main diagonal
in A337388 follow:
\begin{equation}\label{diagrowth}
 T(n,n)\sim\frac{(2+\sqrt n)^{2n+1/2}}{\sqrt{8\pi n}}
 \sim\frac{e^{4\sqrt n-4}n^{n-1/4}}{\sqrt{8\pi}}.
\end{equation}
\end{corollary}
\begin{proof}
For fixed $b=\sqrt k>0$, the same integral as above with exponent
$2n+1$, divided by $b$, is $U(n,k)$; the Gaussian calculation
gives the first formula. This is Kot\v e\v sovec's recorded
A337369 asymptotic.
For the varying parameter put $M_n=b_n+2$ and
$\theta=\sqrt{M_n/(2n)}\,y$ in the even integral.
Since $M_n/n\to0$,
\[
 \left(1-\frac{2(1-\cos\theta)}{M_n}\right)^{2n}
 \longrightarrow e^{-y^2}
\]
for fixed $y$. The base is nonnegative since $b_n\geq2$.
On $[-\pi,\pi]$, $1-\cos\theta\geq2\theta^2/\pi^2$,
so the transformed integrand, extended by zero, is bounded by
$e^{-4y^2/\pi^2}$. Dominated convergence proves the second formula.
Take $b_n=\sqrt n$ for $n\geq4$ and expand
$(2n+1/2)\log(1+2/\sqrt n)=4\sqrt n-4+O(n^{-1/2})$.
This proves \eqref{diagrowth}, including Kot\v e\v sovec's 2020
equivalents. The source's second equivalent includes a factor
$1+19/(3\sqrt n)+199/(18n)$, which tends to $1$; the displayed
$\sim$ is therefore established, without interpreting that factor
as a claimed higher-order error bound.
\end{proof}

\section{Coverage and verification}
References to fields below use A337389 revision 40, A082758 revision 103,
A337390 revision 33, A245926 revision 48, A337419 revision 24,
A337369 revision 30 and A337388 revision 10. ``Proved'' refers to a
statement actually recorded in the source; additional generalizations
are identified separately. No priority assertion is made.
\begingroup\small
\begin{longtable}{@{}p{.43\linewidth}p{.53\linewidth}@{}}
\toprule
Source statement & Coverage in this note\\\midrule
A337389, Bala's final FORMULA conjecture & Reproved from the prior binomial identity: Wildon, 2013; Theorem~\ref{main}, all $n,k\geq0$.\\
A337389, NAME and binomial formula & Proved equal: Theorem~\ref{main}.\\
Polynomial parameter and integral description & Not posed; proved in Theorem~\ref{main} and Corollary~\ref{oddcor}.\\
A337389, recorded $k=0,4$ cases & Reproved as known specializations: Corollary~\ref{consequences}.\\
A082758, April 3, 2022 diagonal conjecture & Reproved; prior proof Bala, 2025: Theorem~\ref{main}, $k=1$.\\
A337390 and A245926, May 2022 diagonal conjectures & Reproved; prior proof Blanco G\'omez, 2026, Section 3.3: $k=2,3$.\\
Equivalent column double/triple sums & Proved by numerator expansion after Theorem~\ref{main}; same prior column proofs.\\
Gauss congruences in those three entries & Proved for all primes: Corollary~\ref{gauss}; all-integer-parameter extension not posed.\\
A082758, $p^{2r}$ conjecture & Prior proof cited: Pan--Sun, 2020, Theorem 1.1; application recorded by Epstein, 2026.\\
A337390 and A245926, $p^{2r}$ conjectures & Reproved from prior proof Blanco G\'omez, 2026, Theorem A and Corollary 3.3; Corollary~\ref{super}.\\
A337390, squared-coefficient conjecture & Proved: Corollary~\ref{squares}; $k=3,5$ also recover the recorded A245926 and A243946 squares.\\
General square formula and varying-column growth & Not posed; proved in Corollaries~\ref{squares} and \ref{growth}.\\
A337389 and A337369, Manyama recurrences & Proved: equation~\eqref{rec} and Corollary~\ref{oddrec}.\\
A337389 and A337369, Kot\v e\v sovec asymptotics & Proved for fixed $k>0$: equation~\eqref{asympt} and Corollary~\ref{growth}.\\
A337419, NAME, binomial formula, recurrence & Proved by $k\mapsto-k$; Corollary~\ref{consequences} and equation~\eqref{rec}.\\
A337388, NAME and two asymptotic equivalents & Proved: Corollaries~\ref{consequences} and \ref{growth}; higher-order error bounds not asserted in that source.\\
A337369, NAME and binomial formula & Proved equal: Corollary~\ref{oddcor}.\\
Columns $k=0,\ldots,5$, named in A337389 & Covered as specializations, including their recorded recurrences and leading growth formulas where present.\\
Other combinatorial interpretations, alternate sums and congruences in column entries & Not covered, except the square and supercongruence formulas explicitly listed above.\\
\bottomrule
\end{longtable}
\endgroup

The Lean declaration
\href{https://github.com/the-omega-institute/trureturing/blob/3378450a2176909e857354ab064dde6ec7679b71/D5/S3/Combinatorics/PowerSeries/BalaA337389Diagonal.lean}{\texttt{BalaA337389Diagonal.result}}
proves the coefficient identity for every natural parameter and degree,
including existence of the normalized pair of square roots over
$\mathbb Q$. Its definitions \texttt{delta}, \texttt{IsColumn}, and
\texttt{phi} are respectively \eqref{defA}'s discriminant and two
normalized square equations, and \eqref{defPhi}.
The polynomial-parameter extension, congruences, square identity,
companion recurrence and asymptotics have the written proofs and cited
theorems given here; these are not additional conclusions of that Lean
declaration. The accompanying exact-rational checker compares finite
coefficient ranges, both radical constructions, recurrences, column sums,
square identities, congruences and recorded initial array entries.

\begin{thebibliography}{99}\small
\bibitem{array} S. Manyama, \emph{Square array defined by a nested
square root}, OEIS A337389, revision 40, May 5, 2022.
Final FORMULA conjecture by P. Bala, May 3, 2022.
\url{https://oeis.org/history/view?seq=A337389&v=40}.
\bibitem{wildon} M. Wildon, answer to \emph{Combinatorial identities},
MathOverflow, November 27, 2013, 20:30 UTC.
\url{https://mathoverflow.net/a/150152}.
\bibitem{gessel2013} I. M. Gessel, answer to \emph{Combinatorial identities},
MathOverflow, November 27, 2013, 16:48 UTC; identification with
terminating Kummer summation.
\url{https://mathoverflow.net/a/150135}.
\bibitem{columnone} P. Bala, diagonal identity and its 2025 proof
addition, and A. R. Epstein, 2026 supercongruence observation,
in OEIS A082758, revision 103, July 18, 2026.
\url{https://oeis.org/history/view?seq=A082758&v=103}.
\bibitem{columntwo} S. Manyama, OEIS A337390, revision 33,
May 5, 2022; conjectures by P. Bala, May 2, 2022.
\url{https://oeis.org/history/view?seq=A337390&v=33}.
\bibitem{columnthree} P. D. Hanna, OEIS A245926, revision 48,
November 10, 2024; conjectures by P. Bala, May 3, 2022.
\url{https://oeis.org/history/view?seq=A245926&v=48}.
\bibitem{blanco} R. Blanco G\'omez, \emph{Exponential transforms and
supercongruences: an order-two transfer theorem, a first-level criterion,
and modular forms behind Bala's prime bounds}, version 1,
October 17, 2026, Theorem A, Corollary 3.3 and Section 3.3.
\url{https://doi.org/10.5281/zenodo.23429108}.
Version 2, same date, retains these results:
\url{https://doi.org/10.5281/zenodo.23430779}.
\bibitem{pansun} H. Pan and Z.-W. Sun, \emph{Supercongruences for
central trinomial coefficients}, arXiv:2012.05121v1, December 9, 2020,
Theorem 1.1. Final version: \emph{Frontiers in Combinatorics and
Number Theory} \textbf{2} (2026), 55--62.
\url{https://arxiv.org/abs/2012.05121v3}.
\bibitem{clausen} NIST Digital Library of Mathematical Functions,
equation 16.12.2 (Clausen's identity).
\url{https://dlmf.nist.gov/16.12.E2}.
\bibitem{signed} S. Manyama, \emph{Signed square array},
OEIS A337419, revision 24, August 29, 2020.
\url{https://oeis.org/history/view?seq=A337419&v=24}.
\bibitem{diagonal} S. Manyama, \emph{Main diagonal of A337389},
OEIS A337388, revision 10, August 31, 2020.
\url{https://oeis.org/history/view?seq=A337388&v=10}.
\bibitem{oddarray} S. Manyama, \emph{Companion square array},
OEIS A337369, revision 30, August 31, 2020.
\url{https://oeis.org/history/view?seq=A337369&v=30}.
\bibitem{columnfive} P. D. Hanna, OEIS A243946, revision 44,
August 23, 2026.
\url{https://oeis.org/history/view?seq=A243946&v=44}.
\end{thebibliography}
\medskip
\noindent\small Copyright \copyright\ 2026 Haobo Ma and Wenlin Zhang.
This note is licensed under Creative Commons Attribution 4.0 International
(CC BY 4.0), \url{https://creativecommons.org/licenses/by/4.0/}.
\end{document}
